\documentclass[a4paper,10pt]{report}

\def\packagepath{../../preambule}   % path du package princial
\usepackage{\packagepath/preambule}    % utilisation du fichier de configuration

\def\level{Première }              % Classe
\def\course{NSI}              % Matière
\def\eval{DS03}

\def\visibleornot{visible}    % visible or invisible
\def\documentpath{./Sources_Latex}
\def\date{Mardi 14 Octobre 2025}
\renewcommand{\arraystretch}{1}  % Ecart dans les tableau

\def\ptsexoA{2}
\def\ptsexoB{2}
\def\ptsexoC{2}
\def\ptsexoD{2}
\def\ptsexoE{2}
\def\ptsexoF{2}
\def\ptsexoG{2}

\def\ptstotal{\ptsexoA+\ptsexoB}

\usepackage{hyperref}

\usetikzlibrary{shapes, arrows, positioning}
\usetikzlibrary{shapes.geometric, arrows, positioning, decorations.pathreplacing}

\tikzstyle{etat} = [draw, rounded corners=15pt, minimum width=2.5cm, minimum height=1.2cm, text centered, font=\bfseries]
\tikzstyle{nouveau} = [etat, fill=blue!40]
\tikzstyle{pret} = [etat, fill=green!60]
\tikzstyle{elu} = [etat, fill=yellow!60]
\tikzstyle{bloque} = [etat, fill=red!60]
\tikzstyle{termine} = [etat, fill=white]
\tikzstyle{fleche} = [thick, ->, >=stealth]

\usepackage{float} % Permet d'utiliser [H]
\usepackage[T1]{fontenc}
\begin{document}

%%%%%%%%%%%%%%%%%%%%%%%%%%%%%%%%%%%%%%%%%%%%%%%%%%ù

%\PageGardeSujetBac[Matiere = Numérique et Sciences Informatiques,
%                  Session = 2025,
%                  AffJour=false,
%                  Duree = {3 heures 30},
%                  DernierePage = \pageref{LastPage},
%                  ModeExamen=false
%                  ]
\renewcommand{\labelitemi}{\textbullet} %pour éviter les tirets dans les "itemize" qui apportent confusion avec le signe moins.
%% Début page de garde style BAC

   %%%%%%%%%%%%%%%%%%%%%%%%%%%%%%%%%%%%%%%%%%%%%%%%%%%%%%%%
%\PageGardeSujetBac[clés]

\pagestyle{DS_FP}          % Feuille de style fancy
\NomPrenomNote{}


\exods{\ptsexoA}


On exécute le script python ci-dessous, que vaut le couple \verb|(a, b)| à la fin de l'exécution?


\begin{CodePythontexAlt}[TaillePolice=\large,Largeur=1\linewidth,Centre,PremLigne=1]{}
def carre(a):
    a = a * a
    return a

a = 2
b = carre(a)
\end{CodePythontexAlt}


    \begin{answer}{1}{2.5}  
    
    \begin{itemize}
        \item         \verb|a| vaut 2. En effet, cette variable n'est pas modifiée à l'extérieur de la fonction. C'est une variable globale qui est passée en paramètre.
        \item \verb|b| vaut ce qui est renvoyé par la fonction \verb|carre| avec comme argument \verb|2|. Dans la fonction \verb|carre|, la variable \verb|a| est une variable locale qui est oubliée dès que la fonction renvoie une valeur (ici \verb|4|.)
    \end{itemize}

    \end{answer}

\medskip

\exods{\ptsexoB}

La fonction \verb|ajoute(n, p)| calcule la somme de tous les entiers entre \verb|n| et \verb|p| inclus. On suppose \verb|n < p|.

Compléter cette fonction

\begin{CodePythontexAlt}[TaillePolice=\large,Largeur=1\linewidth,Centre,PremLigne=1]{}
def ajoute(n, p):
    ''' Ajoute tous les entiers entre n et p inclus
    n et p entiers (int) et n < p
    sortie: entier '''
    somme = 0 
    for i in range(n, p+1):
        somme = somme + i
    return somme 

assert ajoute(2, 4) == 9    # 2 + 3 + 4 = 9
    \end{CodePythontexAlt}


    \medskip

\exods{\ptsexoC}

Ecrire une fonction qui renvoie \verb|True| si \verb|a| est strictement supérieur à \verb|b| et \verb|False| sinon.

\begin{CodePythontexAlt}[TaillePolice=\large,Largeur=1\linewidth,Centre,PremLigne=1]{}
def compare(a, b):
    ''' Renvoie True si a strictement supérieur à b et False sinon
    sortie: bool '''
    if a > b:
        return True
    else:
        return False

def compare2(a, b):
    return a > b

assert compare(1, 5) == False
assert compare(4, 2) == True
    \end{CodePythontexAlt}

 \pagebreak
    \pagestyle{DS_LP}

\exods{\ptsexoD}

Ecrire une fonction \verb|est_present(lettre_cherchee, mot)| qui renvoie \verb|True| si \verb|lettre_cherchee| est dans \verb|mot|, \verb|False| sinon

\begin{CodePythontexAlt}[TaillePolice=\large,Largeur=1\linewidth,Centre,PremLigne=1]{}
def est_present(lettre_cherchee, mot):
    ''' Renvoie True si lettre_cerchee est dans mot et False sinon / sortie: bool '''
    for lettre in mot:
        if lettre == lettre_cherchee:
            return True
    return False

assert est_present('a', 'bonjour') == False
assert est_present('o', 'bonjour') == True
assert est_present('n', 'bonjour') == True

    \end{CodePythontexAlt}

    \medskip

    \exods{\ptsexoE}


Ecrire une fonction \verb|compte_lettre_mot(lettre_cherchee, mot)| qui compte le nombre de \verb|lettre_cherchee| dans \verb|mot|

\begin{CodePythontexAlt}[TaillePolice=\large,Largeur=1\linewidth,Centre,PremLigne=1]{}
def compte_lettre_mot(lettre_cherchee, mot):
    ''' Compte le nombre de lettre_cherchee dans mot / sortie: int '''
    nb_lettre = 0
    for lettre in mot:
        if lettre == lettre_cherchee:
            nb_lettre = nb_lettre + 1
    return nb_lettre

assert compte_lettre_mot('a', 'bonjour') == 0
assert compte_lettre_mot('o', 'bonjour') == 2
assert compte_lettre_mot('n', 'bonjour') == 1
    \end{CodePythontexAlt}

    \medskip

    \exods{\ptsexoF}


Ecrire une fonction \verb|supprime_voyelles(mot)| qui renvoie le \verb|mot| passé en paramètre auquel on a supprimé les voyelles.

\begin{CodePythontexAlt}[TaillePolice=\large,Largeur=1\linewidth,Centre,PremLigne=1]{}

VOYELLES = 'aeiouy'

def supprime_voyelles(mot):
    ''' Renvoie le mot entré en paramètre sans les voyelles / sortie: str '''
    mot_modifie = ''
    for lettre in mot:
        if lettre not in VOYELLES:
            mot_modifie = mot_modifie + lettre
    return mot_modifie

assert supprime_voyelles('bonjour') == 'bnjr'
assert supprime_voyelles('html') == 'html'
assert supprime_voyelles('eau') == ''
    \end{CodePythontexAlt}


\end{document}
